%----------------------------------------------------------------------------------------
%	ABSTRACT PAGE
%----------------------------------------------------------------------------------------
\begin{abstract}
% \addchaptertocentry{\abstractname} % Add the abstract to the table of contents
Modern software systems rely heavily on third-party dependencies and package managers.
This dependence has increased development speed, but it has also expanded the cyber attack surface,
especially in software supply chains. This report studies major cyber threat classes and narrows focus
to software supply chain attacks as a high-impact, ecosystem-level risk.

The report presents: (i) an overview of malware, phishing, denial-of-service, insider threats, and
supply chain attacks; (ii) a case study of the 2024 XZ Utils backdoor incident; (iii) a review of related
research, including MALOSS and dependency confusion analysis; and (iv) a comparative discussion of
existing and proposed mitigations, including eBPF-based monitoring, static analysis, and symbolic
execution assisted triage.

Our key conclusion is that trust itself becomes the primary attack vector in software supply chains.
No single defense is sufficient; robust mitigation requires layered controls across maintainers,
registries, developers, and deployment pipelines. The report closes with practical future directions:
dependency provenance verification, stronger release integrity checks, continuous monitoring, and
faster ecosystem-wide response to newly discovered malicious packages.

\end{abstract}
